\providecommand{\kBKM}{\kappa_{\mathrm{BKM}}}
\providecommand{\kch}{\kappa_{\mathrm{ch}}}
\providecommand{\kcat}{\kappa_{\mathrm{cat}}}
\providecommand{\kfib}{\kappa_{\mathrm{fiber}}}
\providecommand{\Z}{\mathbb{Z}}
\providecommand{\cO}{\mathcal{O}}
\providecommand{\Hup}{\mathbb{H}_{2}}
\providecommand{\wt}{\operatorname{wt}}
\providecommand{\Ltwo}{L^{\mathrm{CHL}}_{2}}
\providecommand{\DeltaTwo}{\Delta_{4}^{(2)}}
\providecommand{\DeltaTwoGC}{\Delta_{2}^{\mathrm{GC}}}
\providecommand{\PhiTwoScal}{\Phi_{2}^{\mathrm{scal}}}
\providecommand{\PhiPrim}{\Phi^{\mathrm{prim}}}

\section*{$N=2$ primitive CHL denominator}
\label{sec:N2-CHL-automorphic-denominator}

\paragraph{Primitive CHL row.}
Let $S$ be a K3 surface, let $g_2$ be a symplectic involution of
Mathieu class $2A$ with frame shape $1^8 2^8$, and let
$t_2\in E[2]$ be a non-zero translation of an elliptic curve.  The
free diagonal quotient
\[
  X_2=(S\times E)/\langle(g_2,t_2)\rangle
\]
is a smooth Calabi--Yau threefold: $t_2$ makes the diagonal action free,
and $g_2$ preserves $H^{2,0}(S)$.  Write
$\Gamma_0^+(2)\subset\mathrm{Sp}_4(\mathbb Q)$ for the Fricke extension
of the level-two genus-two congruence group in the Gritsenko--Nikulin
normalisation, equivalently the orthogonal modular group attached to
\[
  \Ltwo = U\oplus U(2)\oplus\langle -2\rangle ,
  \qquad \operatorname{sig}(\Ltwo)=(2,3).
\]
The CHL-averaged primitive Borcherds denominator is
\[
  \DeltaTwo\in M_4(\Gamma_0^+(2),\chi_2),
\]
the $N=2$ member of the Gritsenko--Nikulin tower
$\Delta_{k(N)}^{(N)}$ with $k(2)=4$.  Its divisor is the Humbert surface
$H_1(\Gamma_0^+(2))\subset \Hup/\Gamma_0^+(2)$ in the Gritsenko--Nikulin
product chamber, and its character is $\chi_2$ (Gritsenko--Nikulin
1995, Part~II, Theorem~2.1).  The scalar square convention is separate:
\[
  \PhiTwoScal=(\DeltaTwo)^2,
  \qquad \wt(\PhiTwoScal)=8.
\]
It records the scalar reciprocal-denominator convention alone; the
primitive automorphic denominator, Hall construction, and positive-half
construction are separate data.

\paragraph{Borcherds input constant.}
Let
\[
  \varphi^{\mathrm{CHL}}_2(\tau,z)
  =Z^{(2A)}_{K3}(\tau,z)
\]
be the weight-$0$, index-$1$ weak Jacobi input on $\Gamma_0(2)$ in the
CHL-averaged convention.  It is normalised by
\[
  \varphi^{\mathrm{CHL}}_2(\tau,z)\big|_{q^0}
  =
  \zeta^{-1}+8+\zeta ,
  \qquad \zeta=e^{2\pi iz}.
\]
Thus the Borcherds constant is
\[
  c^{\mathrm{CHL}}_2(0)
  =[q^0\zeta^0]\varphi^{\mathrm{CHL}}_2=8.
\]
The same value follows from the theta-component supertrace:
\[
  T_{H^*}(g_2)=16,\qquad A_2=4,\qquad
  c^{\mathrm{CHL}}_2(0)=T_{H^*}(g_2)-2A_2=16-8=8.
\]
Here $T_{H^*}(g_2)$ is the Lefschetz trace on the Mukai cohomology and
$A_2$ is the odd theta-component coefficient of the $2A$ twined genus.
Borcherds' singular theta-lift weight formula (Borcherds 1998,
Theorem~13.3), applied in the Gritsenko--Nikulin product chamber, gives
\[
  \wt(\DeltaTwo)=\frac{c^{\mathrm{CHL}}_2(0)}{2}=4,
  \qquad
  \kBKM(\DeltaTwo)=4.
\]
This is the $N=2$ entry of the primitive CHL tuple
\[
  (c_N(0))_{N=(1,2,3,4,6)}=(10,8,6,4,2),
  \qquad
  (\kBKM(\PhiPrim_N))_{N=(1,2,3,4,6)}=(5,4,3,2,1).
\]
Here $\PhiPrim_N$ denotes the primitive CHL denominator; at $N=2$,
$\PhiPrim_2=\DeltaTwo$.

\paragraph{Diagonal-divisor row.}
The Gritsenko--Cl\'ery diagonal-divisor form denoted by $\Delta_2$ in
that catalogue is written here as $\DeltaTwoGC$.  Its row data are
\[
  (t,N_{\mathrm{lvl}};k)=(2,1;2),
  \qquad
  \DeltaTwoGC\in M_2(\Gamma_2,\chi_4),
  \qquad
  \operatorname{div}(\DeltaTwoGC)=\mathcal H_1
\]
on $\Hup/\Gamma_2$, with $\mathcal H_1$ the diagonal divisor of
multiplicity one.  Its product lattice is
\[
  M_2=U(8)\oplus\langle 2\rangle,
  \qquad \operatorname{sig}(M_2)=(2,1),
\]
with exponential $q^n r^\ell s^{2m}$.  Its Jacobi input is the
Gritsenko--Nikulin weak Jacobi form $\phi_{0,2}$ in the singly-twined
Eichler--Zagier/Gritsenko--Cl\'ery normalisation, with
$c^{\mathrm{GC}}_2(0)=4$, hence $\wt(\DeltaTwoGC)=2$.  These data belong
to the Gritsenko--Cl\'ery catalogue, whose row system is
triple-indexed by $(t,N_{\mathrm{lvl}};k)$.  They are separate from the
primitive CHL data $c^{\mathrm{CHL}}_2(0)=8$ and $\DeltaTwo$ of weight
$4$.

\paragraph{Product normalisation.}
For $Z=(\tau,z,\sigma)$ and
$(q,\zeta,p)=(e^{2\pi i\tau},e^{2\pi iz},e^{2\pi i\sigma})$, the
automorphic product has the form
\[
  \DeltaTwo(Z)
  =
  q^A\zeta^B p^C
  \prod_{\substack{(n,\ell,m)\in\Z^3\\(n,\ell,m)>0}}
  \left(1-q^n\zeta^\ell p^m\right)^{
    c^{\mathrm{CHL}}_2(nm,\ell)} .
\]
The positivity convention is the cusp convention of the
Gritsenko--Nikulin product:
\[
  (n,\ell,m)>0
  \Longleftrightarrow
  m>0\ \text{or}\ (m=0,\ n>0)\ \text{or}\ (m=n=0,\ \ell<0).
\]
The monomial $q^A\zeta^B p^C$ is the Weyl-vector factor in this chamber.
The coefficient $c^{\mathrm{CHL}}_2(0)=8$ determines the automorphic
weight.  The coefficient system $c^{\mathrm{CHL}}_2(n,\ell)$ determines
the product exponents and the signed imaginary-root multiplicities of
the target denominator algebra.

\paragraph{Recognition data.}
The automorphic product supplies target-side denominator data.  A BKM
recognition theorem additionally specifies
\[
\begin{gathered}
  \Ltwo=U\oplus U(2)\oplus\langle -2\rangle,\\
  \text{the real simple roots and Weyl chamber,}\\
  \text{the signed imaginary-root multiplicities from }
  c^{\mathrm{CHL}}_2(n,\ell),\\
  \text{and the Weyl--Kac--Borcherds identity with this chamber.}
\end{gathered}
\]
Gritsenko--Nikulin supply this target-side denominator datum for the
CHL paramodular product.  The weight computation proves
$\kBKM(\DeltaTwo)=4$.  A compact Hall or CY-to-chiral realisation is a
separate theorem: it must supply CoHA multiplication, orientation data,
parity, radical quotient, PBW comparison, and bracket compatibility with
the protected BPS source.  Since $X_2$ is a Calabi--Yau threefold, the
native output of $\Phi_3$ is $E_1$-chiral.  Braided $E_2$ structure
belongs to a centre, double, or module layer after these additional data
are constructed.

\paragraph{$K3\times E$ invariants.}
For the quotient $X_2$, finite \'etale descent gives
$H^{0,q}(X_2)=H^{0,q}(S\times E)^{\langle(g_2,t_2)\rangle}$.
The involution is symplectic on $S$ and the translation is trivial on
$H^{0,\bullet}(E)$, hence the dimensions are $1,1,1,1$ in degrees
$0,1,2,3$.  Thus
\[
  \kcat(X_2)=\chi(\cO_{X_2})=0,
  \qquad
  \kch(\Phi_3(X_2))=\sum_q(-1)^q h^{0,q}(X_2)=0.
\]
The fixed-locus Euler number $\chi(S^{g_2})=8$ equals
$c^{\mathrm{CHL}}_2(0)$ in this order-two convention because it enters
the K3 twining calculation.  It serves as twining input, separate from
the $K3\times E$ $\kappa_\bullet$ invariants.  The BKM invariant is
\[
  \kBKM(\DeltaTwo)=c^{\mathrm{CHL}}_2(0)/2=4.
\]
For $Y=K3\times E$ the standard values remain
\[
  \kcat(Y)=0,\qquad
  \kch(\Phi_3(Y))=0,\qquad
  \kappa_{\mathrm{ch}}^{\mathrm{Heis}}(Y)=3,\qquad
  \kBKM(\Delta_5)=5,\qquad
  \kfib(Y)=24,
\]
each produced by its own construction.

\paragraph{Scalar square comparison.}
The reduced scalar theory uses its own chamber and variable convention.
In the Bryan--Oberdieck order-two reduced curve-counting normalisation
(Bryan--Oberdieck 2018, Theorem~0.1), the reciprocal scalar has the
square-level denominator
\[
  Z^{\mathrm{red}}_{\mathrm{DT}}(X_2)
  = -\,(\PhiTwoScal)^{-1}
  = -\,(\DeltaTwo)^{-2}
\]
after that scalar convention is fixed.  This statement supplies none of
the Weyl sign, Hall source, parity fixture, or Lie bracket.  A BPS Lie
algebra comparison must match the compact graded multiplicities with
$c^{\mathrm{CHL}}_2(nm,\ell)$ and prove compatibility with the
Kontsevich--Soibelman product.  The positive half, the Hall--Drinfeld
double, and any Oberdieck--Pandharipande scalar comparison are separate
inputs.

\paragraph{The $N=1$ anchor.}
The $N=1$ row fixes the primitive normalisation:
\[
  \PhiPrim_1=\Delta_5,\qquad
  \kBKM(\Delta_5)=10/2=5.
\]
The scalar Igusa convention uses $\Phi_{10}^{\mathrm{scal}}=(\Delta_5)^2$.
This $N=1$ square leaves the $N=2$ Gritsenko--Cl\'ery diagonal-divisor
row separate from the primitive CHL denominator.

\paragraph{Primary references.}
Mukai 1988, Proposition~1.2 and Table~\S2, gives the fixed-locus data
for the symplectic involution.  Eguchi--Hikami 2011, Table~1, gives the
$2A$ theta-component correction $A_2=4$.  Gritsenko--Nikulin 1995,
Part~II, Theorem~2.1, constructs the paramodular forms
$\Delta_{k(N)}^{(N)}$ with $k(2)=4$ in the CHL-averaged convention.
Gritsenko 1999 gives the additive-lift normalisation of the same
integer-weight ladder.  Borcherds 1998, Theorem~13.3, gives the weight
formula $c_0(0)/2$ for the singular theta lift.  Gritsenko--Cl\'ery
2008, Theorems~1.4 and~3.1, govern the separate diagonal-divisor
catalogue, including the row $(t,N_{\mathrm{lvl}};k)=(2,1;2)$.
Bryan--Oberdieck 2018, Theorem~0.1, gives the reduced CHL scalar in the
square-level convention.
